We describe here a concrete instantiation of our PPB scheme using a lossy variant of the ElGamal encryption scheme \cite{C:ElGamal84}.
% \input{figures/folklore_lossy_elgamal}
The lossy variant, considered cryptographic ``folklore,'' uses an algorithm $\egpphgen$ that takes as input a security parameter $\secparam$ and outputs a group description with \emph{two} generators $G = (\group, \egprime, \ggen, \hgen)$, where the discrete log of $\hgen$ with respect to $\ggen$ is unknown.
The scheme generates lossy keys by running the traditional ElGamal key generation algorithm with the alternate generator $h$.
As the first element of a ciphertext remains $g^r$, computational lossiness follows from DDH.

We begin by proving that this folklore scheme has the properties we require of it.
Next, we give a transformation from a weak collision-free to a weak robust PKE.
We then use these tools to describe concrete instantiations for $\pker$ and $\pkem$.
Finally, we give some optimizations to these instantiations that compress both the size of the escrow value and the public key.

\subsection{Lossy ElGamal}

A lossy PKE based on ElGamal can be seen in \cref{fig:folklore_lossyelg}.
\begin{figure}[h]
	\begin{pcvstack}[boxed, center, space=1em]
		\begin{pchstack}[center, space=1em]
			\procedure[linenumbering]{$\setup(\secparam)$}{
				\pcreturn \egpphgen(\secparam)
			}
			\procedure[linenumbering]{$\kgen(\pkepp)$}{
				\pcparse (\group, \egprime, \ggen, \hgen) = \pkepp \\
				x \sample \Z_{\egprime} \\
				X \leftarrow \ggen^x \\
				pk \leftarrow(\egprime, \ggen, X) \\
				sk \leftarrow(\egprime, \ggen, x) \\
				\pcreturn (sk, pk)
			}
			\procedure[linenumbering]{$\enc(\pkepp, pk, M)$}{
				\pcparse (\group, \egprime, \ggen, \hgen) = \pkepp \\
				y \sample \Z_{\egprime} \\
				Y \leftarrow \ggen^y \\
				T \leftarrow X^y \\
				W \leftarrow T M \\
				\pcreturn (Y, W)
			}
		\end{pchstack}
		\begin{pchstack}[center, space=1em]
			\procedure[linenumbering]{$\dec(\pkepp, sk, ct)$}{
				\pcparse (\group, \egprime, \ggen, \hgen) = \pkepp \\
				(Y, W) \leftarrow ct  \\
				T \leftarrow Y^x \\
				M \leftarrow W T^{-1} \\
				\pcreturn M
			}
			\procedure[linenumbering]{$\kgenlossy(\pkepp)$}{
				\pcparse (\group, \egprime, \ggen, \hgen) = \pkepp \\
				r \sample \bin^{\lambda} \\
				X \leftarrow h^r \\
				pk \leftarrow (\egprime, \ggen, X) \\
				\pcreturn pk
			}
		\end{pchstack}
	\end{pcvstack}
	\caption{A lossy PKE adapted from the ElGamal encryption scheme, $\lelgamal$. }
	\label{fig:folklore_lossyelg}
\end{figure}


\begin{theorem}
	If the Decisional-Diffie-Hellman assumption holds relative to $\egpphgen$, then the scheme shown in \cref{fig:folklore_lossyelg} satisfies computational lossiness.
\end{theorem}

\begin{proof}
	We prove security using a sequence of hybrids.

	\begin{itemize}
		\item $\newhyb$: This hybrid is identical to $\complossygame_0$.
		\item $\nexthyb$: In this hybrid $\ct = (g^y, (h^{ry}\msg_0))$ is replaced with $(g^y, (h^{rz}\msg_0))$ for $z \sample \mathbb{Z}_p$.
		      We then have that $\prevhyb \cindist \currhyb$, as otherwise we could build a reduction $\bdv$ to DDH as follows.
		      $\bdv$ receives $(g, h, g^y, c)$ from the challenger, and $(r, m_0, m_1)$ from $\adv$.
		      $\bdv$ sends $(g^y, c^r\msg_0)$ to $\adv$, and outputs what $\adv$ outputs.
		      In the case that $c = h^y$, $\adv$'s view is exactly as it is in $\prevhyb$, while when $c = h^z$, $\adv$'s view is exactly as it is in $\currhyb$.
		\item $\nexthyb$: In this hybrid $\ct$ is replaced with $(g^y, (h^{rz}\msg_1))$.
		      As $z$ is a uniformly random element of $\mathbb{Z}_p$ the message is fully masked, and so $\prevhyb \pindist \currhyb$.
		\item $\nexthyb$: In this hybrid $\ct$ is replaced with $(g^y, (h^{ry}\msg_1))$.
		      We have that $\prevhyb \cindist \currhyb$ from DDH as before.
		      Note that $\currhyb$ is identical to $\complossygame_1$.
	\end{itemize}
	We therefore have that $\hyb_0 \cindist \currhyb$, and so the scheme satisfies computational lossiness.

\end{proof}

\begin{theorem}
	The scheme shown in \cref{fig:folklore_lossyelg} satisfies $\mathcal{U}_\group$-strong-lossy-ambiguity.
\end{theorem}

\begin{proof}
	We prove security using a sequence of hybrids.

	\begin{itemize}
		\item $\newhyb$: This game is identical to $\stronglossyambgame_0$.
		\item $\nexthyb$: In this hybrid $\ct = (g^y, pk_0^y\cdot \msg)$ is replaced with $\ct = (g^y, pk_1^y\cdot \msg)$.
		      As $m$ is sampled uniformly at random from $\group$, we have that $\prevhyb \pindist \currhyb$.
	\end{itemize}

	Then, as $\currhyb$ is identical to $\stronglossyambgame_1$, we have that the scheme satisfies $\mathcal{U}_\group$-strong-lossy-ambiguity.
\end{proof}

\subsection{A Weak Collision-Free to Robust Transform}

We give here a generic transformation from a weak collision-free PKE to a robust one.
The transform can be seen in \cref{fig:weak_robust_transform}.

\begin{figure}[ht]
	\begin{pcvstack}[boxed, center, space=1em]
		\begin{pchstack}[center, space=1em]
			\procedure[linenumbering]{$\setup'(\secparam)$}{
				\flag \sample \msgspace \\
				\pcreturn (\setup(\secparam), \flag)
			}
			\procedure[linenumbering]{$\kgen'((\pkepp, \flag))$}{
				(sk_0, pk_0) \leftarrow \kgen(\secparam) \\
				(sk_1, pk_1) \leftarrow \kgen(\secparam) \\
				\pairsk = (sk_0, sk_1) \\
				\pairpk = (pk_0, pk_1) \\
				\pcreturn (\pairsk, \pairpk)
			}
			\procedure[linenumbering]{$\kgenlossy'((\pkepp, \flag))$}{
				pk_0 \leftarrow \kgenlossy(\secparam) \\
				pk_1 \leftarrow \kgenlossy(\secparam) \\
				\pairpk = (pk_0, pk_1) \\
				\pcreturn \pairpk
			}
		\end{pchstack}
		\begin{pchstack}[center, space=1em]
			\procedure[linenumbering]{$\enc'((\pkepp, \flag), pk, \msg)$}{
				(pk_0, pk_1) \leftarrow pk \\
				\ct_0 \leftarrow \enc(pk_0, \msg) \\
				\ct_1 \leftarrow \enc(pk_1, \flag) \\
				\pcreturn (\ct_0, \ct_1)
			}
			\procedure[linenumbering]{$\dec'((\pkepp, \flag), sk, \ct)$}{
				(sk_0, sk_1) \leftarrow sk \\
				(\ct_0, \ct_1) \leftarrow \ct \\
				\pcif \dec(sk_1, \ct_1) \neq \flag \\
				\t \pcreturn \bot  \\
				\pcelse \\
				\t \pcreturn \dec(sk_0, \ct_0)
			}
		\end{pchstack}
	\end{pcvstack}
	\caption{A transform from a collision-free $\pke$ to a weakly-robust scheme $\pke'$.}
	\label{fig:weak_robust_transform}
\end{figure}


\begin{theorem}
	If $\pke$ satisfies weak collision-freeness, then the scheme $\pke'$ shown in \cref{fig:weak_robust_transform} satisfies weak robustness.
\end{theorem}

\begin{proof}
	Assume there exists an adversary $\adv$ such that $\advantage{\robust}{\pke', \adv}$ is non-negligible.
	We can then construct a reduction $\bdv$ to the collision-freeness of $\pke$ as follows.

	$\bdv$ receives $\{pk_i\}_{i\in[n]}$, computes $\{pk_i'\}_{i\in [n]} \leftarrow \kgen(\secparam)$, runs $\adv(\secparam, \{(pk'_i,$ $ pk_i)\}_{i\in[n]})$, receives $(j, k, m, r = (r_0, r_1))$, and outputs $(j, k, \phi, r_1)$.
	As $\adv$'s advantage is non-negligible, it must be the case that $\dec(sk_k, \enc(pk_j, \phi; r_1')) = \phi$ with non-negligible probability.
	Therefore $\bdv$'s advantage is non-negligible, which is a contradiction, and so the scheme satisfies weak robustness.
\end{proof}

\begin{theorem}
	If $\pke$ satisfies weak-lossy-ambiguity, then the scheme $\pke'$ shown in \cref{fig:weak_robust_transform} satisfies weak-lossy-ambiguity.
\end{theorem}

\begin{proof}
	We prove security using a sequence of hybrids.

	\begin{itemize}
		\item $\newhyb$: This hybrid is identical to $\wlambgame_0$.
		\item $\nexthyb$: In this hybrid $\ct_0 \leftarrow \enc(pk_{0, 0}, m)$ is replaced with $\enc(pk_{1, 0}, m)$.
		      We then have that $\prevhyb \cindist \currhyb$ as otherwise we could build a reduction $\bdv$ to the \wla{} of $\pke$ as follows.
		      $\bdv$ receives $(r_0 = (r_{0, 0}, r_{0, 1}), r_1 = (r_{1, 0}, r_{1, 1}), m)$ from $\adv$, and sends $(r_{0, 0}, r_{1, 0}, m)$ to the challenger, receiving $\ct^*$.
		      $\bdv$ then sends $(\ct^*, \enc(\kgenlossy(;r_{0, 1}), \phi))$ to $\adv$, and outputs whatever $\adv$ outputs.
		\item $\nexthyb$: In this hybrid $\ct_1$ is replaced by $\enc(pk_{1, 1}, \phi)$.
		      We again have that $\prevhyb \cindist \currhyb$ by the \wla{} of $\pke$.
	\end{itemize}

	Then, as $\currhyb$ is identical to $\wlambgame_1$, we have that $\hyb_0 \cindist \currhyb$ and $\pke'$ satisfies \wla{}.
\end{proof}

\subsection{Instantiation}
\label{sec:instantiation}

$\pker$ can be instantiated as the result of applying the transform show in \cref{fig:weak_robust_transform} to the lossy ElGamal variant $\lelgamal$ shown in \cref{fig:folklore_lossyelg}.
Next, we need to ensure that the message space of $\pkem$ is the same as the ciphertext space of $\pker$.
Each $\pker$ ciphertext is $4$ group elements, and so we can have $\pkem$ be four copies of $\lelgamal$ run in parallel.


\subsection{Optimizations}
\label{sec:ppb-optimizations}
We describe here optimizations that reduce the public-key and escrow size of the concrete PPB scheme described in \cref{sec:instantiation} while preserving security.

\subsubsection{Randomness Reuse.}
As written, a $\pker$ public-key would consist of two sub-keys $pk_0 = g^{x_0}$, $pk_1 = g^{x_1}$, and a ciphertext would consist of two ElGamal encryptions of the form
$((g^y, g^{x_0y}\cdot m), (g^z, g^{x_1z} \cdot \phi))$.
We can compress this ciphertext via randomness re-use.
ElGamal encryption allows an encryptor to re-use the same $g^r$ value across encryptions under different public-keys.
This was first shown to be secure in \cite{PKC:Kurosawa02}, and later tested securely under a more stringent security definition in \cite{PKC:BelBolSta03}.
We can leverage this to use the same $g^{r_{\rev}}$ value across \emph{all} reveal ciphertexts, as each is under a different key.
$\pker$ ciphertexts now take the form:
$(g^{x_0r_{\rev}}\cdot m, g^{x_1r_{\rev}} \cdot \phi)$,
with $g^{r_{\rev}}$ included separately in the escrow value.

\subsubsection{Shared Public Keys.}
After the above optimization match ciphertexts $\matct$ must consist of \emph{two} ElGamal encryptions.
For these, we can use the same $(pk_0, pk_1)$ present in the $\pker$ public key.
We can additionally use the same randomness-reuse trick as before, using the same random value $r_{\mat}$ for all match ciphertexts.
A match ciphertext $\matct$ encrypting a pad $\otpk = (\otpk_0, \otpk_1)$ then has the form:
$(g^{x_0r_{\mat}} \cdot \otpk_0, g^{x_1r_{\mat}} \cdot \otpk_1)$,
with $r_{\mat}$ included separately in the escrow value.
With this optimization the scheme's public-key only needs to include two ElGamal public-keys in each cell, rather than four as before.


\subsection{Proofs}
We prove here that the instantiation from \cref{sec:instantiation} satsifies lossy correlation resistance.

\begin{theorem}
	The instantiations of $\pker$, $\pkem$ described in \cref{sec:instantiation}, satisfy lossy-correlation-resistance if the discrete logarithm problem is hard for $\egpphgen$.
\end{theorem}

\begin{proof}
	We begin by recalling from \cref{sec:ppb-optimizations} that $\pker$ and $\pkem$ are both made up of two copies of an underlying ElGamal PKE scheme over a group $\group$ run in parallel.
	We will abuse notation and use the same variables above to describe only the \emph{second} copy of the scheme (i.e.\ the one in the $\pker$ construction that encrypts $\flag$).
	For readability we will also refrain from listing the $g^y$ (i.e.\ randomness) portion of ElGamal ciphertexts.

	Assume that there exists an adversary $\adv$ such that $\advantage{\badcorgame}{\pker, \pkem,  \adv}$ is non-negligible.
	We can then construct a reduction $\bdv$ to the discrete log problem for $\egpphgen$.
	$\bdv$ receives $(g, h)$, samples $(\rkgr_0, \rkgr_1, \mkgr_0, \mkgr_1) \sample \Z$, invokes $\adv$ on $((g, h), h^{\mkgr_1}, g^{\mkgr_0},$ $g^{\rkgr_0}, g^{\rkgr_1})$, and receives $(\mencr, \rencr, \okgrA, \okgrB, \okgrC, \okgrD)$.
	$\bdv$ then outputs $\okgrA + \okgrB + \okgrD + \mkgir\mencr + \rkgrA\rencr - \rkgrB\rencr - \okgrC - \mkglr\mencr$.

	As $\adv$ has non-negligible advantage, we must have that with non-negligible probability:
	\begin{align*}
		\pker.\dec(\revsk, (\otpk_1^{-1} \cdot \pkem.\dec(\matsk, \matct))^{-1} \cdot (\otpk_2 \cdot \pker.\enc(\revpk, \otpk^*; \rencr))) & \\
		\neq \bot                                                                                                                          &
	\end{align*}

	\vspace{-2em}
	\begin{align*}
		(\otpk_1^{-1} \cdot \pkem.\dec(\matsk, \matct))^{-1} \cdot (\otpk_2 \cdot \pker.\enc(\revpk, \otpk^*; \rencr))                      & = g^{\rkgrB\rencr} \tag{Definition of $\pker$}              \\
		\left(\otpk_1^{-1} \cdot \frac{h^{\mkglr\mencr}\otpk_3}{g^{\mkgir\mencr}}\right)^{-1} \cdot (\otpk_2 \cdot g^{\rkgrA\rencr}\otpk^*) & = g^{\rkgrB\rencr} \tag{Definitions of $\pker$ and $\pkem$} \\
		\frac{\otpk_1\otpk_2\otpk^* g^{\mkgir\mencr}g^{\rkgrA\rencr}}{\otpk_3h^{\mkglr\mencr}}                                              & = g^{\rkgrB\rencr}                                          \\
		\frac{\otpk_1\otpk_2\otpk^* g^{\mkgir\mencr}g^{\rkgrA\rencr}}{\otpk_3g^{\rkgrB\rencr}}                                              & = h^{\mkglr\mencr}                                          \\
		g^{\okgrA + \okgrB + \okgrD + \mkgir\mencr + \rkgrA\rencr - \rkgrB\rencr - \okgrC}                                                  & = h^{\mkglr\mencr}                                          \\
		\okgrA + \okgrB + \okgrD + \mkgir\mencr + \rkgrA\rencr - \rkgrB\rencr - \okgrC - \mkglr\mencr                                       & = s
	\end{align*}
	Where $s$ is the discrete log of $h$ with respect to $g$.

	Therefore $\bdv$ succeeds with non-negligible probability, which contradicts our assumption that discrete log is hard in $\group$, and so $(\pker, \pkem)$ satisfy lossy-correlation-resistance.
\end{proof}

